%-------------------------
% Resume in LaTeX
% Two-column layout: section labels on the left, content on the right,
% hairline rules between sections. Black-and-white, ATS-friendly.
%-------------------------

\documentclass[letterpaper,10pt]{article}

\usepackage{fontspec}
\setmainfont{Ubuntu}[
    Path=fonts/,
    Extension=.ttf,
    UprightFont=Ubuntu-Regular,
    BoldFont=Ubuntu-Bold,
    ItalicFont=Ubuntu-Italic,
    BoldItalicFont=Ubuntu-BoldItalic,
    Numbers=Monospaced
]
\usepackage{enumitem}
\usepackage[dvipsnames,table]{xcolor}
\usepackage[hidelinks]{hyperref}
\usepackage{geometry}
\usepackage{changepage}
\usepackage{graphicx}
\usepackage{needspace}

\geometry{
    letterpaper,
    left=0.7in,
    right=0.7in,
    top=0.7in,
    bottom=0.65in
}

\definecolor{accent}{HTML}{000000}
\definecolor{subtle}{HTML}{3F3F3F}

\hypersetup{
    colorlinks=true,
    urlcolor=accent,
    linkcolor=accent
}

\pagestyle{empty}
\setlength{\parindent}{0pt}
\setlist[itemize]{leftmargin=1.2em, topsep=4pt, itemsep=2pt, parsep=0pt, label=\textbullet}

%-------- Layout machinery --------
% Each section is a resumesection environment: a hairline rule (except before
% the first section), then content indented by \labelw. The section label is
% hung into the left margin on the first line of the first entry.
\newlength{\labelw}
\setlength{\labelw}{1.45in}
\newlength{\nextgap}
\newcounter{secn}
\newcommand{\sectionlabel}{}

% Starts an entry: vertical gap, then the pending section label (if any)
\newcommand{\entrystart}{%
  \par\vspace{\nextgap}\global\nextgap=9pt\relax\noindent
  \ifx\sectionlabel\empty\else
    \kern-\labelw
    \parbox[t]{\dimexpr\labelw-1em\relax}{\raggedright\small\bfseries\addfontfeatures{LetterSpace=0.5}\sectionlabel}%
    \kern1em
    \global\let\sectionlabel\empty
  \fi
}

% #1=LABEL (write in capitals)
\newenvironment{resumesection}[1]{%
  \par\needspace{6\baselineskip}%
  \stepcounter{secn}%
  \ifnum\value{secn}>1
    \par\vspace{11pt}\noindent\hspace*{\labelw}\rule{\dimexpr\linewidth-\labelw\relax}{0.5pt}\par\vspace{11pt}%
  \fi
  \gdef\sectionlabel{#1}\global\nextgap=0pt\relax
  \begin{adjustwidth}{\labelw}{0pt}%
}{%
  \end{adjustwidth}\par
}

%-------- Entry macros --------
% #1=icon path (unused)  #2=company  #3=location  #4=role  #5=dates
\newcommand{\jobheading}[5]{%
  \entrystart\textbf{#4\ \textbar\ #5}\\[2pt]
  \textcolor{subtle}{#2, #3}\par
  \vspace{-1pt}
}

% #1=title  #2=date
\newcommand{\plainheading}[2]{%
  \entrystart\textbf{#1\ \textbar\ #2}\par
  \global\nextgap=6pt\relax
}

% #1=icon path (unused)  #2=title  #3=date (optional, may be empty)
\newcommand{\iconentry}[3]{%
  \entrystart #2\if\relax\detokenize{#3}\relax\else\ \textbar\ \textbf{#3}\fi\par
  \global\nextgap=5pt\relax
}

% #1=title  #2=date
\newcommand{\projheading}[2]{%
  \entrystart\textbf{#1\ \textbar\ #2}\par
}

% #1=icon path (unused)  #2=school/degree  #3=detail  #4=dates (optional, may be empty)
\newcommand{\eduheading}[4]{%
  \entrystart\textbf{#2\if\relax\detokenize{#4}\relax\else\ \textbar\ #4\fi}\\[2pt]
  \textcolor{subtle}{#3}\par
}

% Sub-role heading within a single job entry (e.g. dual-hat responsibilities)
\newcommand{\subrole}[1]{%
  \par\vspace{4pt}\noindent{\small\itshape\color{subtle}#1}\par\vspace{-1pt}%
}

% Highlight tag row beneath job/project title --- appears outside bullet list
% #1=label (e.g. Stack, Focus)  #2=content
\newcommand{\stackline}[2]{%
  \par\vspace{4pt}\noindent
  \small\textbf{#1:}~\textcolor{subtle}{#2}\par
  \vspace{1pt}
}

% Skills row: #1=category  #2=items
\newcommand{\skillrow}[2]{%
  \entrystart\small\textbf{#1:} #2\par
  \global\nextgap=5pt\relax
}

\begin{document}

%-------- HEADER --------
\noindent{\fontsize{27}{31}\selectfont\bfseries Sajid Pervez}\hfill{\large AI Cyber Security \& Controls Engineer}\par
\vspace{22pt}

%-------- CONTACT --------
\begin{resumesection}{CONTACT}
\entrystart
\begin{minipage}[t]{0.5\linewidth}\small
  \textbf{Phone}: 0469 366 692\\
  \textbf{Email}: \href{mailto:sajid@sajidpervez.com}{sajid@sajidpervez.com}
\end{minipage}%
\begin{minipage}[t]{0.5\linewidth}\small
  \textbf{LinkedIn}: \href{https://linkedin.com/in/sajidpervez}{linkedin.com/in/sajidpervez}\\
  \textbf{Location}: Melbourne, VIC
\end{minipage}
\end{resumesection}

%-------- SUMMARY --------
\begin{resumesection}{SUMMARY}
\entrystart\small AI cyber security and controls leader with 15+ years securing software, cloud, and AI/ML systems. Hands-on experience defining AI security requirements, assessing AI risk, leading threat modelling, and implementing governance controls for production AI, RAG, and agentic systems on AWS. Brings deep NIST, ISM, ASD Essential Eight, secure-SDLC, vulnerability-management, and stakeholder-engagement experience. Prior NV1 security clearance (inactive).\par
\end{resumesection}

%-------- EXPERIENCE --------
\begin{resumesection}{PROFESSIONAL EXPERIENCE}

\jobheading{icons/company.png}{MYOB}{Melbourne, VIC}{Principal AI Security Engineer}{May 2025 -- Present}
\subrole{AI security governance, controls, and risk leadership for AWS-hosted products}
\begin{itemize}
    \small
    \item Lead AI security architecture reviews and threat modelling for multiple production AI products, identifying model-evasion, adversarial-attack, privacy-leakage, and prompt-injection risks early in delivery.
    \item Define AI security requirements and secure-by-design controls for prompt and input handling, response guardrails, and policy-as-code governance.
    \item Secure RAG pipelines end to end through retrieval-integrity validation, embedding security, indirect prompt-injection prevention, and data-provenance controls.
    \item Write security requirements for MCP servers and agentic workflows, applying the OWASP LLM Top 10 and MITRE ATLAS frameworks.
    \item Conduct break-testing and red-teaming exercises to surface vulnerabilities before production deployment; implement observability and monitoring controls for continuous improvement.
    \item Partner with Cyber Security, AI Engineering, and business stakeholders to establish technical direction and embed AI security controls in engineering workflows.
\end{itemize}
\stackline{Focus}{AI governance and controls, AI risk management, agentic workflows, RAG security, prompt injection, MCP security, guardrails, OWASP LLM Top 10, MITRE ATLAS, AWS}

\jobheading{icons/company.png}{MYOB}{Melbourne, VIC}{Principal Product Security Engineer}{Jun 2022 -- Apr 2025}
\subrole{Product security controls, secure software delivery, and cloud-native application protection}
\begin{itemize}
    \small
    \item Led product security initiatives through solution-design reviews and threat modelling using STRIDE, OWASP Top 10, and CWE/SANS Top 25.
    \item Rolled out CNAPP tooling (Wiz and Orca) for continuous cloud-security posture monitoring and real-time AWS infrastructure visibility.
    \item Created vulnerability dashboards for leadership, improving risk prioritisation, accountability, and visibility across AI and cloud environments.
    \item Implemented automated DevSecOps pipelines integrating SAST, SCA, and DAST, enabling repeatable vulnerability detection throughout the secure software development lifecycle.
    \item Led architecture and design discussions for Agile delivery, contributing risk-based guidance, threat-model estimates, and continuous security-improvement cycles across multidisciplinary teams.
\end{itemize}
\stackline{Stack}{AWS, cloud-native security, Wiz, Orca, SAST, SCA, DAST, CI/CD, secure SDLC, STRIDE, OWASP, CWE, NIST CSF, ISM, ASD Essential Eight}

\jobheading{icons/company.png}{NBN Australia}{Australia}{Application Security Lead}{Sep 2020 -- Jun 2022}
\begin{itemize}
    \small
    \item Uplifted enterprise application security across complex multi-team delivery environments using OWASP SAMM and NIST CSF.
    \item Established security-maturity baselines aligned to the ISM and ISO/IEC 27001.
    \item Rolled out SAST and SCA capabilities with automated incident reporting to improve vulnerability detection across high-risk applications.
    \item Implemented secure-code review processes and a security-champion program, embedding secure engineering practices across delivery teams.
\end{itemize}
\stackline{Stack}{OWASP SAMM, NIST CSF, ISM, ISO/IEC 27001, SAST, SCA, vulnerability management, secure code review}

\jobheading{icons/company.png}{AGL Energy}{Australia}{Application Security Engineer}{Aug 2019 -- Sep 2020}
\begin{itemize}
    \small
    \item Designed security-scanning architecture within Azure DevOps CI/CD pipelines using SAST, SCA, and DAST tooling, embedding automated security gates.
    \item Rolled out a vulnerability-management platform and SDL operational model, improving incident-response workflows and mean time to remediation.
\end{itemize}
\stackline{Stack}{Azure DevOps, SAST, SCA, DAST, CI/CD security, vulnerability management, SDL}

\plainheading{Application Security Consultant, Telstra}{2018 -- 2019}
\plainheading{Senior TS Advisor / PCI QSA, UL Transaction Security}{2015 -- 2018}
\plainheading{PKI \& Strong Authentication Developer, Various}{2005 -- 2015}
\end{resumesection}

%-------- EDUCATION --------
\begin{resumesection}{EDUCATION}
\eduheading{icons/edu.png}{Master of Computer Science}{Postgraduate degree}{}
\eduheading{icons/edu.png}{Bachelor of Computer Science}{Undergraduate degree}{}
\end{resumesection}

%-------- CERTIFICATIONS --------
\begin{resumesection}{CERTIFICATES}
\iconentry{icons/cert.png}{\textbf{CISSP} \textcolor{subtle}{(ISC)\textsuperscript{2} \quad$\cdot$\quad CSSLP (ISC)\textsuperscript{2}}}{}
\iconentry{icons/cert.png}{\textbf{AWS Security Engineer} \quad$\cdot$\quad \textbf{PCI QSA} \quad$\cdot$\quad \textbf{PA-DSS QSA}}{}
\end{resumesection}

%-------- SKILLS --------
\begin{resumesection}{SKILLS}
\skillrow{AI Governance \& Controls}{AI security requirements, AI risk assessment, AI governance controls, policy-as-code, threat modelling, responsible deployment, monitoring and assurance}
\skillrow{AI/ML Security}{LLM guardrails, prompt-injection mitigation, data-poisoning prevention, model protection, adversarial-attack defence, privacy-leakage controls, RAG security, MCP and agent security, OWASP LLM Top 10, MITRE ATLAS}
\skillrow{AWS Security}{Security Hub, GuardDuty, IAM, KMS, WAF, Macie, CNAPP (Wiz, Orca)}
\skillrow{Compliance}{ISO/IEC 27001, NIST CSF, PCI DSS, OWASP ASVS, ISM, ASD Essential Eight}
\skillrow{DevSecOps}{SAST, SCA, DAST, static and dynamic security testing, CI/CD pipeline security, GitOps, secure MLOps, automated security gates}
\skillrow{Security Architecture}{Security architecture reviews, secure-by-design requirements, threat modelling, security design patterns, risk-based guidance, cloud-native application security}
\skillrow{Application Security}{OWASP SAMM, STRIDE, OWASP Top 10, CWE/SANS Top 25, secure code review, secure SDLC, SDL}
\skillrow{Programming}{Python, JavaScript/Node.js, Bash, C\#/.NET, Java}
\skillrow{Delivery}{Agile, Scrum, incident response management, vulnerability assessment, red teaming}
\end{resumesection}

\end{document}
